\section{Background}
\label{sec:background}

\paragraph{KV blocks are buffer-pool pages} A transformer's attention state for a token prefix is
stored as fixed-size blocks keyed by a hash of the prefix, so two requests sharing a prefix share
blocks. Serving systems keep these blocks across a hierarchy of GPU memory, host DRAM, remote memory and
storage, and decide which tier holds which block~\cite{qin2024mooncake,gao2024cachedattention,
F5_liu2025lmcache,zhang2026polarkv}. From a database perspective this is buffer management with a new page
type~\cite{pan2025dbperspective}: pages are immutable once written, a page's first access is a
compulsory miss (its KV is computed, not read), and a miss on a cached page costs a transfer rather
than a recomputation. Multi-tier buffer managers have long treated the movement between tiers as a
device-priced policy decision~\cite{canim2010ssdbufferpool,zhou2021spitfire,wu2019anna}; Spitfire, for
one, migrates lazily and with tunable probabilities, and may bypass a tier, because a move has a price.
Most KV-cache studies do not price moves (below); the ski-rental and cascade-aware rules of
\Cref{sec:online} are the simplest members of that lazy-migration family.

\paragraph{Replacement and admission} LRU-K~\cite{oneil1993lruk} ranks pages by the time of their
$K$-th most recent reference and states the case this paper turns on: ``The only time the choice is
ambiguous is when more than one page has [infinite backward $K$-distance]. In this case, a subsidiary
policy may be used.'' Blocks with no usable history are that case, and on KV traces they are most
blocks. 2Q~\cite{johnson1994twoq} and its descendants admit new pages to a probationary queue;
S3-FIFO~\cite{yang2023s3fifo} and SIEVE~\cite{zhang2024sieve} demote one-hit wonders quickly;
TinyLFU~\cite{einziger2017tinylfu} admits a new page only if its estimated frequency beats the victim's.
ARC~\cite{S7_megiddo2003arc}, LeCaR~\cite{S8_vietri2018lecar}, CACHEUS~\cite{S9_rodriguez2021cacheus} and LRB~\cite{S6_song2020lrb} adapt or learn the replacement
decision, and learning-augmented heuristics tune a simple policy's parameters
instead~\cite{xia2026lah}. In KV serving, HotPrefix~\cite{li2025hotprefix} admits hot prefixes to host
memory and promotes them to the GPU, and OrbitFlow~\cite{ma2026orbitflow} prices GPU--host transfers
inside one request's plan.

\begin{table}[t]
\centering\footnotesize\setlength{\tabcolsep}{3pt}
\begin{tabular}{@{}p{2.2cm}p{2.1cm}p{1.75cm}p{1.75cm}@{}}
\toprule
study & when it decides & bound reported & moves priced \\
\midrule
CachedAttention~\cite{gao2024cachedattention} & online; warm-up on 10K turns, evaluate on the next 42K & none & storage, not moves \\
Mooncake~\cite{qin2024mooncake} & online, LRU store & infinite cache & in scheduling, not eviction \\
KVCache in the wild~\cite{wang2025kvcachewild} & online & infinite cache & ``negligible'' \\
Fancy eviction~\cite{liu2026fancyeviction} & simulated online & Belady, compute-aware & compute, not moves \\
PolarKV~\cite{zhang2026polarkv} & online LRU, background offload & none & budget, not moves \\
Baleen~\cite{S12_wong2024baleen} & trained on day 1, tested after & episode OPT & flash writes \\
\midrule
this work & split instant and online replay & references with stated tie rules & measured; five ledgers \\
\bottomrule
\end{tabular}
\caption{How cache studies close to ours evaluate, from each paper's evaluation section. All evaluate online
policies, whose own admission rule orders new blocks; none uses a split-time reference, the setting where
the tie rule of \Cref{sec:ties} decides the result.}
\label{tab:protocols}
\end{table}

\paragraph{How KV-cache placement is evaluated} \Cref{tab:protocols} lists the closest studies. They
warm a cache on a prefix of the trace and measure on the rest, or bound a policy by an infinite cache
or by Belady's offline optimum. Their policies are online, so a policy's own admission rule orders the
blocks it has not seen, and we do not claim that their reported numbers move with our tie rule. The tie
rule decides a different kind of result: a split-time placement analysis, or a reference for how much
of an oracle's advantage a predictor could reach, which is the evaluation this paper first ran
(\Cref{sec:ties}). Movement between tiers is assumed free, hidden behind computation, or priced
as provisioning, with Baleen's flash-write budget the exception. A concurrent study of fourteen eviction
policies on production prefix-cache traces~\cite{liu2026fancyeviction} finds plain LRU within 0.4 points
of the best online policy at a 1\,TiB pool and Belady far above all of them. It overlaps ours in three
observations: block lifetimes are bounded by sessions, early frequency predicts whether a block is reused
at all, and quick demotion helps only where one-hit traffic dominates. It prices recomputation, not moves
between tiers, and does not vary how the measurement starts. What this paper adds is the tie rule's effect
on a reference, a ladder that prices what each part of an oracle's knowledge is worth, and the price and
ledger of a move.
